\documentclass[aspectratio=169,11pt]{beamer}

% --- Theme ---
\usetheme{metropolis}
\usepackage{appendixnumberbeamer}

% --- Packages ---
\usepackage{fontspec}
\usepackage{minted}
\setminted{fontsize=\small,breaklines=true}
\usepackage{booktabs}
\usepackage{tikz}
\usepackage{fontawesome5}

% --- Colors ---
\definecolor{healthgreen}{RGB}{34,139,34}
\definecolor{alertred}{RGB}{220,53,69}
\setbeamercolor{alerted text}{fg=alertred}

% --- Metadata ---
\title{Module 5: Visualization for health data}
\subtitle{Data analysis with Python for health specialists}
\author{Ya\'{e} Ulrich Gaba}
\institute{AIRINA Labs}
\date{2026}

\begin{document}

% ============================================================
\maketitle

% ============================================================
\begin{frame}{Principles of health data visualization}
Before writing any plotting code, ask:
\begin{enumerate}
  \item \textbf{What is the message?} ``Malaria incidence is declining'' --- not ``here is data.''
  \item \textbf{Who is the audience?} Clinicians need different plots than policymakers.
  \item \textbf{What comparison matters?} Between groups? Over time? Against a threshold?
\end{enumerate}

\vspace{6pt}
Rules for \alert{every} health figure:
\begin{itemize}
  \item Label axes with units (``Systolic BP (mmHg)'', not ``BP'')
  \item Use colorblind-friendly palettes
  \item No 3D charts, no pie charts, no dual y-axes
  \item Include sample sizes in titles or annotations
\end{itemize}
\end{frame}

% ============================================================
\begin{frame}[fragile]{Setup: matplotlib and seaborn}
\begin{minted}{python}
import matplotlib.pyplot as plt
import seaborn as sns

# Clean style for all plots
sns.set_theme(style="whitegrid", font_scale=1.1)
plt.rcParams["figure.figsize"] = (8, 5)
plt.rcParams["figure.dpi"] = 100

# Colorblind-friendly palette (Bang Wong, 2011)
cb_palette = ["#0072B2", "#D55E00", "#009E73",
              "#CC79A7", "#F0E442", "#56B4E9"]
sns.set_palette(cb_palette)
\end{minted}

These six colors are distinguishable by people with the most common forms of color vision deficiency.
\end{frame}

% ============================================================
\begin{frame}[fragile]{Histograms: distribution of a variable}
\begin{minted}{python}
fig, ax = plt.subplots(figsize=(10, 5))
ax.hist(all_glucose, bins=40, edgecolor="white", alpha=0.7)
ax.axvline(100, color="orange", linestyle="--",
           label="Normal threshold (100 mg/dL)")
ax.axvline(126, color="red", linestyle="--",
           label="Diabetic threshold (126 mg/dL)")
ax.set_xlabel("Fasting glucose (mg/dL)")
ax.set_ylabel("Number of patients")
ax.set_title("Distribution of fasting glucose (n=1,000)")
ax.legend()
\end{minted}

\vspace{4pt}
Add \alert{clinical threshold lines} to give immediate context.
\end{frame}

% ============================================================
\begin{frame}[fragile]{Box plots: group comparisons}
\begin{minted}{python}
fig, ax = plt.subplots(figsize=(10, 6))
sns.boxplot(data=clinical, x="age_group", y="bmi",
            hue="sex", order=["18-39", "40-59", "60+"])
ax.axhline(25, color="orange", linestyle="--",
           label="Overweight")
ax.axhline(30, color="red", linestyle="--",
           label="Obese")
ax.set_ylabel("BMI (kg/m\u00b2)")
ax.set_title(f"BMI by sex and age group (n={n})")
ax.legend()
\end{minted}

\begin{alertblock}{Limitation}
Box plots hide bimodal distributions. Consider \texttt{sns.violinplot()} or overlay \texttt{sns.stripplot()} for small samples.
\end{alertblock}
\end{frame}

% ============================================================
\begin{frame}[fragile]{Bar charts: counts and rates}
\begin{minted}{python}
# Horizontal bar chart: 15 lowest life expectancy countries
bottom15 = recent.nsmallest(15, "lifeExp")

fig, ax = plt.subplots(figsize=(10, 7))
ax.barh(bottom15["country"], bottom15["lifeExp"],
        color="#D55E00")
ax.set_xlabel("Life expectancy (years)")
ax.set_title("15 countries with lowest life expectancy, 2007")
ax.invert_yaxis()
for i, v in enumerate(bottom15["lifeExp"]):
    ax.text(v + 0.5, i, f"{v:.1f}", va="center")
\end{minted}

Use \textbf{horizontal bars} when you have many categories to avoid rotated labels.
\end{frame}

% ============================================================
\begin{frame}[fragile]{Line plots: epidemic curves}
\begin{minted}{python}
fig, ax = plt.subplots(figsize=(12, 5))
ax.bar(epi["date"], epi["new_cases"], alpha=0.4,
       label="Daily cases")
ax.plot(epi["date"], epi["rolling_7d"], color="#D55E00",
        linewidth=2, label="7-day average")
ax.set_ylabel("New cases per day")
ax.set_title("Epidemic curve: daily new cases")
ax.legend()
\end{minted}

\begin{exampleblock}{Clinical context}
The epi curve is the \alert{single most important visualization} during an outbreak. Its shape reveals the transmission pattern: point-source (sharp peak), continuous (plateau), or person-to-person (waves).
\end{exampleblock}
\end{frame}

% ============================================================
\begin{frame}[fragile]{Scatter plots: relationships}
\begin{minted}{python}
# The classic Gapminder bubble chart
fig, ax = plt.subplots(figsize=(10, 7))
for continent in recent["continent"].unique():
    sub = recent[recent["continent"] == continent]
    ax.scatter(sub["gdpPercap"], sub["lifeExp"],
               s=sub["pop"] / 1e6, alpha=0.6,
               label=continent)
ax.set_xscale("log")
ax.set_xlabel("GDP per capita (log scale, USD)")
ax.set_ylabel("Life expectancy (years)")
ax.set_title("Wealth vs Health, 2007")
ax.legend()
\end{minted}

Add regression lines with \texttt{sns.regplot()} and report the Pearson $r$ in the title.
\end{frame}

% ============================================================
\begin{frame}[fragile]{Kaplan-Meier survival curves}
\begin{minted}{python}
from lifelines import KaplanMeierFitter

kmf = KaplanMeierFitter()
kmf.fit(treatment_time, event_observed=treatment_event,
        label="Treatment (n=100)")
kmf.plot_survival_function(ax=ax, ci_show=True)
\end{minted}

\vspace{4pt}
Survival analysis tracks \textbf{time until an event} (death, relapse). KM curves handle \alert{censored} observations --- patients still alive at study end.

\vspace{4pt}
\textbf{Install:} \texttt{pip install lifelines}
\end{frame}

% ============================================================
\begin{frame}[fragile]{Heatmaps: correlations and co-occurrence}
\begin{minted}{python}
corr = health.corr()
mask = np.triu(np.ones_like(corr, dtype=bool))

fig, ax = plt.subplots(figsize=(9, 8))
sns.heatmap(corr, mask=mask, annot=True, fmt=".2f",
            cmap="RdBu_r", center=0, vmin=-1, vmax=1,
            square=True, ax=ax)
ax.set_title("Correlation matrix of health indicators")
\end{minted}

\vspace{4pt}
Use \texttt{"RdBu\_r"} (diverging) for correlations. Use \texttt{"YlOrRd"} (sequential) for disease burden.
\end{frame}

% ============================================================
\begin{frame}[fragile]{Multi-panel figures: the health dashboard}
\begin{minted}{python}
fig, axes = plt.subplots(2, 2, figsize=(14, 10))
# Panel A: Box plot by continent
# Panel B: GDP vs Life Expectancy scatter
# Panel C: Life expectancy trends over time
# Panel D: Bottom 10 countries bar chart
plt.tight_layout()
plt.savefig("dashboard.png", dpi=150, bbox_inches="tight")
\end{minted}

\vspace{4pt}
Combine related plots into a single figure for reports and publications. Save with \texttt{dpi=300} for print quality.
\end{frame}

% ============================================================
\begin{frame}{What you can do after this module}
\begin{enumerate}
  \item Create histograms with clinical threshold lines
  \item Build box plots and bar charts for group comparisons
  \item Plot epidemic curves with 7-day rolling averages
  \item Generate scatter plots with regression lines and $r$ values
  \item Produce Kaplan-Meier survival curves
  \item Combine everything into multi-panel dashboards
\end{enumerate}

\vspace{12pt}
\centering
\Large \alert{Next: Module 6 --- Hypothesis testing}
\end{frame}

% ============================================================
\begin{frame}{Exercises (Hour 3)}
\begin{enumerate}
  \item \textbf{Modify dashboard:} Rebuild the Gapminder dashboard for 2002 data
  \item \textbf{Clinical dashboard:} 4 panels --- BMI histogram, BP box plot by sex, HbA1c vs.\ age scatter, diagnosis frequency bar chart
  \item \textbf{Survival:} Kaplan-Meier plot for 3 treatment groups with median survival annotations
  \item \textbf{Mini-project:} Build a 4-panel health dashboard using Gapminder 2007 data
\end{enumerate}
\end{frame}

\end{document}
